\section{Q-Learning 系统实践与监控}

\subsection{监控指标与数据质量}

\begin{knowledgebox}{关键观测指标}
\begin{itemize}
    \item Replay buffer 覆盖率：衡量最近 $N$ 个转移是否覆盖常见状态-动作对。
    \item TD error 分布：监控 TD error 的标准差与最大值，可提前感知数值爆炸。
    \item Action distribution drift：行为策略的采样频率是否偏向少数动作，防止训练走向狭窄策略。
\end{itemize}
\end{knowledgebox}

为了做到稳定部署，工程团队必须把上述指标暴露到训练仪表盘，并结合基准数据制定警戒线。摄取的经验应被打标签（如挑战性、稀疏 reward、环境变化）以便后续分析。

\begin{importantbox}{Replay Buffer 覆盖与分布警示}
如果 buffer 恰好被早期成功经验撑满，后续探索的样本无法进入训练，这会导致 agent 过早收敛。推荐设置 \texttt{min\_size} 和 \texttt{max\_size} 并定期用采样接口检查时间戳，必要时使用 \texttt{prioritized replay} 或 \texttt{rank-based} 重新激励采样。
\end{importantbox}

\subsection{经验回放策略与数据重平衡}

\begin{knowledgebox}{Experience Replay 重平衡技巧}
\begin{itemize}
    \item Prioritized Replay：以 $|\delta|^\alpha$ 作为采样权重，$\delta$ 为 TD error，$\alpha$ 控制偏好强度。
    \item Importance Sampling Correction：对优先样本做比例补偿，防止 policy bias。
    \item Diversity Buffer：保留来自不同子任务或者模拟 seed 的经验片段，避免训练陷入 narrow distribution。
\end{itemize}
\end{knowledgebox}

直接将 replay buffer 当作一个黑匣子是不够的，建议在每次训练迭代后记录新旧样本进入比例及其平均 reward。若新样本比例低于 30\%，说明探索效率下降，需要加大学习率或 epsilon。

\begin{warningbox}{Replay buffer 过度依赖旧数据}
随着训练推进，旧数据可能已经过时，尤其在非静态环境中。继续训练时要关注 buffer 中平均 timestep，若超过 50\% 的数据来自过去十分钟，可启用循环清洗或将优先级重新置为高值，避免旧策略主导梯度。
\end{warningbox}

\subsection{本章小结}

本章强调 Q-Learning 系统建设中的可观测性与 replay buffer 的健康度，任何一次大规模复训都应在数据收集、采样分布与 TD error 这三个维度设定 guardrail。

